problem:  
Let $(M^n, g)$ be a closed Riemannian manifold whose curvature tensor is **harmonic**, i.e. $\nabla^* \mathrm{Riem} = 0$. Equivalently, the divergence of the Riemann curvature tensor vanishes, but the metric is *not* assumed to be Einstein. Classical results show that locally symmetric spaces automatically satisfy $\nabla \mathrm{Riem} = 0$, but the harmonic curvature condition $\nabla^*\mathrm{Riem} = 0$ is weaker, encompassing certain Kähler and conformally Einstein examples.  

**Question:**  
Does every compact Riemannian manifold $(M,g)$ satisfying $\nabla^*\mathrm{Riem}=0$ and not locally symmetric admit a canonical geometric decomposition  
\[
(M,g) \approx (M_1 \times \cdots \times M_k, g_1 \oplus \cdots \oplus g_k)
\]
up to finite cover, such that each $(M_i, g_i)$ is either (a) Einstein, or (b) conformally Einstein or locally a warped product with Einstein base?  
Equivalently, is the class of compact harmonic-curvature metrics finitely generated under the operations of product, warping over Einstein spaces, and conformal deformation?  

Why is it a "good" problem:  
This formulation corrects the redundancy in combining “harmonic curvature” and “Einstein” (since Einstein implies harmonic curvature). It instead probes the *structural origin* of compact harmonic-curvature metrics, an open classification problem lying between the theory of symmetric spaces and the broader landscape of Riemannian manifolds with special curvature identities.  

The problem targets whether the analytic condition $\nabla^*\mathrm{Riem} = 0$ secretly imposes an underlying geometric composition law—akin to the de Rham, Cheeger–Gromoll, or Berger–Simons decompositions—but in a weaker curvature-harmonic framework.  
A resolution would bridge global Riemannian geometry, conformal geometry, and the analytic theory of curvature operators, revealing whether harmonic-curvature manifolds sit just on the boundary of symmetry or form a genuinely new flexible geometric class.  

The statement is concise, logically sound, and of a level appropriate for research mathematics since it aims to extend deep structural rigidity theorems into a less constrained analytic setting. Its potential to unify diverse examples under a single geometric decomposition principle makes it a genuinely “good” and forward-looking problem.